\documentclass[12pt]{article}
\usepackage{amsmath, amssymb}
\usepackage{geometry}
\usepackage[ruled,vlined]{algorithm2e}
\usepackage{booktabs}
\geometry{margin=1in}

\title{Homework: Branch and Bound}
\author{Algorithm Design and Analysis}
\date{}

\begin{document}
\maketitle

\section{Problem Statement}

Describe the Branch-and-Bound (B\&B) framework and apply it to the 0/1 Knapsack problem.
Demonstrate the search process on a concrete instance.

\section{Branch and Bound Framework}

Branch-and-Bound is a systematic enumeration strategy that:
\begin{enumerate}
  \item \textbf{Branches}: splits the feasible region into sub-problems (children).
  \item \textbf{Bounds}: computes an upper bound (for maximization) for each sub-problem.
  \item \textbf{Prunes}: discards sub-problems whose upper bound cannot exceed the current best solution.
\end{enumerate}

\begin{algorithm}[H]
\SetAlgoLined
\KwIn{Root node of search tree}
\KwOut{Optimal solution}
Initialize priority queue $PQ$ with root node\;
$best \leftarrow -\infty$\;
\While{$PQ$ is not empty}{
  $node \leftarrow$ dequeue node with highest upper bound from $PQ$\;
  \If{$node.\text{upperBound} \leq best$}{\textbf{continue} (prune)}
  \If{$node$ is a leaf}{$best \leftarrow \max(best, node.\text{value})$; \textbf{continue}}
  \ForEach{child of $node$}{
    Compute child.upperBound\;
    \If{child.upperBound $> best$}{enqueue child into $PQ$}
  }
}
\Return $best$\;
\caption{Generic Branch-and-Bound (Best-First)}
\end{algorithm}

\section{Application: 0/1 Knapsack}

\subsection*{Upper Bound Computation}

Relax the 0/1 constraint to allow fractional items (fractional knapsack bound).
Sort remaining items by value density $v_i/w_i$ and greedily fill the knapsack.

\subsection*{Example}

$n=4$, $W=10$. Items (sorted by $v/w$ descending):

\begin{center}
\begin{tabular}{ccccc}
\toprule
Item & Weight & Value & $v/w$ \\
\midrule
1 & 2 & 10 & 5.0 \\
2 & 4 & 16 & 4.0 \\
3 & 6 & 18 & 3.0 \\
4 & 3 & 6  & 2.0 \\
\bottomrule
\end{tabular}
\end{center}

Root node: capacity 10, value 0.
Upper bound = $10 + 16 + (10-2-4)/6 \times 18 = 10+16+12 = 38$.

\textbf{Branch on item 1:}
\begin{itemize}
  \item Include 1: capacity $= 8$, value $= 10$, UB = $10 + 16 + 18 \times (2/6) = 32$.
  \item Exclude 1: capacity $= 10$, value $= 0$, UB = $16 + 18 + 6 \times (1/3) = 36$.
\end{itemize}

Continue expanding until leaf; optimal solution found: items 1, 2, 3 (value 44, weight 12 — infeasible).
Correct optimal: items 1, 2, 4 (weight $2+4+3=9 \leq 10$, value $10+16+6=32$).

\section{Complexity}

Worst case: exponential $O(2^n)$. In practice, effective pruning dramatically reduces nodes explored.

\end{document}
